\subsection{Distribution and conventional-range behaviour of heritability estimates}

The distribution of empirical SNP-heritability estimates depended strongly on reporting scale. On the native observed scale, 759 evaluable configuration-level mean estimates were available across the nine binary phenotypes, of which 121 (15.9\%) were below zero, 638 (84.1\%) were within the conventional interval $0\leq h^2\leq1$, and none exceeded one. Following liability-scale harmonisation, 729 binary configuration--phenotype estimates were evaluable; 118 (16.2\%) were below zero, 467 (64.1\%) were within 0--1, and 144 (19.8\%) exceeded one. For BMI, the only quantitative empirical phenotype, 11 of 85 estimates (12.9\%) were below zero, 73 (85.9\%) were within 0--1, and one (1.2\%) exceeded one. Values outside the conventional range were retained rather than truncated because unconstrained estimators can return negative or greater-than-one values when sampling variability or model instability is substantial.

\subsection{Estimator uncertainty and cross-fold stability}

Estimator-specific uncertainty and cross-fold stability were analysed separately because they quantify different sources of variation. Estimator-specific SEs were available for all evaluable empirical configuration--phenotype units from nine of the ten method groups; LDpred2 did not provide estimator-specific SE in the parsed benchmark output. On the observed binary scale, median estimator-specific SE was 0.874 for GCTA, 0.624 for GEMMA1, and 0.620 for DPR+GEMMA. Following liability-scale transformation, the corresponding median SEs were 2.433, 1.834, and 1.829, respectively. Cross-fold SD was treated separately as a descriptive measure of sensitivity to the overlapping training partitions: on the liability scale, median cross-fold SD was 0.694 for GCTA, 0.758 for GEMMA1, and 0.733 for DPR+GEMMA. Several summary-statistic workflows showed little or no cross-fold variation because the heritability-estimation inputs were largely fixed across folds; near-zero cross-fold SD for these workflows therefore should not be interpreted as evidence of superior inferential precision.

The controlled simulations demonstrated that finite target sample size materially amplified instability. On the binary liability scale, median cross-fold SD decreased from 0.225 at $N=750$ to 0.055 at $N=5{,}000$ for GCTA ($q=1.92\times10^{-47}$), from 0.248 to 0.063 for GEMMA1 ($q=7.93\times10^{-181}$), and from 0.245 to 0.062 for DPR+GEMMA ($q=7.47\times10^{-177}$). Similar reductions were observed for continuous traits. These results demonstrate a substantial finite-sample contribution to empirical instability. However, configuration-dependent effects and known-truth accuracy differences persisted in the larger simulations, indicating that workflow sensitivity cannot be attributed solely to statistical uncertainty caused by limited sample size.